\documentclass[11pt,a4paper]{letter}

\usepackage[utf8]{inputenc}
\usepackage[T1]{fontenc}
\usepackage[margin=2.5cm]{geometry}
%\usepackage{microtype}
\usepackage{hyperref}

\signature{Linwei Tao \\ on behalf of all authors}
\address{School of Computer Science \\ University of Sydney \\ Sydney, NSW 2006, Australia \\ linwei.tao@sydney.edu.au}
\date{\today}

\begin{document}

\begin{letter}{
  Editor-in-Chief \\
  IEEE Transactions on Pattern Analysis and Machine Intelligence
}

\opening{Dear Editor-in-Chief,}

We are pleased to submit our manuscript titled
\textbf{``Confidence Calibration under Ambiguous Ground Truth''}
for consideration for publication in
\textit{IEEE Transactions on Pattern Analysis and Machine Intelligence}.

\textbf{Summary of contributions.}
Standard post-hoc calibration assumes a unique ground-truth label per input and fits
calibrators against majority-voted labels.
We show that this practice is fundamentally flawed under label ambiguity:
Temperature Scaling is theoretically biased toward temperatures that underestimate
annotator uncertainty, with the true-label miscalibration gap growing monotonically
with annotation entropy.
To address this, we propose a family of ambiguity-aware post-hoc calibrators that
optimise proper scoring rules against the full annotator distribution and require no
model retraining.
Our methods span three practical annotation regimes:
(i)~Dirichlet-Soft leverages full annotator distributions and reduces true-label ECE
by 55--87\% relative to Temperature Scaling across four benchmarks;
(ii)~Monte Carlo Temperature Scaling with a single annotation per example (MCTS~$S{=}1$)
matches full-distribution calibration within 0.6~pp ECE across all eight settings,
demonstrating that pre-aggregated label distributions are unnecessary;
and (iii)~Label-Smooth Temperature Scaling (LS-TS) operates with voted labels alone
and reduces ECE by 9--77\% without any annotator data.

\textbf{Relevance to TPAMI.}
This work addresses a fundamental gap in the evaluation and improvement of model
reliability for vision and language systems operating under genuine label ambiguity.
We provide both theoretical analysis and large-scale empirical validation across
vision (CIFAR-10H, ISIC~2019, DermaMNIST) and language (ChaosNLI) benchmarks,
covering eight architecture configurations.
We believe this contribution is well-suited to the TPAMI audience, which spans
pattern recognition, computer vision, and machine learning.

\textbf{Originality.}
This manuscript has not been published previously and is not under consideration
for publication elsewhere.
All authors have approved the manuscript and agree with the submission to TPAMI.
The work does not raise any ethical concerns; all datasets used are publicly available
or derived from clinician-reported statistics in the published literature.

We suggest the manuscript be classified as a \textbf{Regular Paper}.

We look forward to receiving the reviewers' feedback and thank you for your time
and consideration.

\closing{Sincerely,}

\end{letter}
\end{document}
